\documentclass[11pt,letterpaper]{article}
\usepackage{fontspec}
\usepackage{tocloft}
\usepackage{multicol}
\usepackage{nicefrac}
\usepackage[left=1.3in,right=1.3in,top=1in,bottom=1in]{geometry}
\setmainfont[Scale=1.1,AutoFakeBold=1.5,AutoFakeSlant=0.3]{Doves Type}

% Prevent page breaks within list items
\widowpenalty=10000
\clubpenalty=10000
\interlinepenalty=500

\title{Napa Cabbage Kimchi (Mak-Kimchi)}
\author{}
\date{}

\begin{document}

\maketitle
\thispagestyle{empty}
\section*{Ingredients}
\setlength{\columnsep}{20pt}
\begin{multicols}{2}
\noindent
    Napa cabbage \dotfill 5--6~lb. \\
    Diamond Crystal kosher salt \dotfill \nicefrac{3}{4}~cup \\
    Water \dotfill 2~cups \\
    Mochiko (sweet rice flour) \dotfill 3~Tbsp. \\
    Sugar \dotfill 1\nicefrac{1}{2}~Tbsp. \\
    Garlic, minced \dotfill \nicefrac{1}{2}~cup \\
    \columnbreak
    Fresh ginger, grated \dotfill 1~Tbsp. \\
    Yellow onion, grated \dotfill \nicefrac{1}{2}~medium \\
    Fuji apple, grated \dotfill \nicefrac{1}{2}~medium \\
    Fish sauce \dotfill \nicefrac{1}{2}~cup \\
    Gochugaru, coarse \dotfill 1\nicefrac{1}{2}~cups \\
    Daikon radish \dotfill 12~oz. \\
    Carrot, medium \dotfill 1 \\
    Scallions \dotfill 8 \\
\end{multicols}

\section*{Directions}

\noindent
Cut 5--6~lb. \textbf{napa cabbage} into 2~inch pieces; set aside 4 large outer leaves (these go on top of the packed kimchi, inside the jars) ---
Mince \textbf{garlic} to \nicefrac{1}{2}~cup; grate 1~Tbsp. \textbf{ginger}, \nicefrac{1}{2}~medium \textbf{yellow onion}, and \nicefrac{1}{2}~medium \textbf{Fuji apple} ---
Cut 12~oz. \textbf{daikon} and 1 medium \textbf{carrot} into thin matchsticks; cut 8~\textbf{scallions} into 1~inch lengths ---
Wash four 1-quart wide-mouth mason jars in hot soapy water; rinse and dry ---
Have gloves ready for mixing

\begin{enumerate}
    \item Put the 5--6~lb. cut \textbf{napa cabbage} in a 12-qt. mixing bowl. Sprinkle \nicefrac{3}{4}~cup \textbf{Diamond Crystal kosher salt} over it, working the salt into the thick white ribs. Toss every \textit{30~minutes} for \textit{1\nicefrac{1}{2}--2~hours}, until the white ribs bend with a little spring, the pile has slumped, liquid pools in the bowl, and a bite tastes quite salty. Continue salting in \textit{15~minute} stretches until the ribs bend.

    \item While the cabbage salts, make the porridge. In a small saucepan, whisk 3~Tbsp. \textbf{Mochiko} into 2~cups cold \textbf{water} until smooth. Cook over medium heat, stirring, \textit{5--8~minutes}, until the mixture is glossy, slightly translucent, and thick enough to coat a spoon, with no dry flour lumps. Stir in 1\nicefrac{1}{2}~Tbsp. \textbf{sugar} and cook \textit{1~minute} more. Take it off the heat. Cool until it is no longer warm to the touch, \textit{20--30~minutes}. A warm porridge cooks the garlic and thins the paste.

    \item Tip the salted cabbage into a colander. Rinse under cold water \textit{2--3~times}, lifting and turning so the salt in the ribs washes out. Taste a thick piece and continue rinsing until it is pleasantly salty, not harsh. Drain \textit{20--30~minutes}, then squeeze gently in handfuls until no streams of water run, but the cabbage still feels moist. Return it to the 12-qt. mixing bowl. Cabbage wrung bone-dry will not make enough brine to fizz and stay covered.

    \item Mix the paste, then the kimchi:
    \begin{enumerate}
        \item Add the cooled porridge, \nicefrac{1}{2}~cup minced \textbf{garlic}, 1~Tbsp. grated \textbf{ginger}, the grated \nicefrac{1}{2}~\textbf{onion}, the grated \nicefrac{1}{2}~\textbf{Fuji apple}, \nicefrac{1}{2}~cup \textbf{fish sauce}, and 1\nicefrac{1}{2}~cups coarse \textbf{gochugaru}. With gloved hands, mix until the seasonings form a thick red paste and coat the cabbage. Taste a coated piece: it should be saltier and spicier than you want the finished kimchi. Add a pinch of salt until it tastes vivid, or a splash more \textbf{fish sauce} until it tastes savory.
        \item Add the 12~oz. matchstick \textbf{daikon}, the matchstick \textbf{carrot}, and the 8 cut \textbf{scallions}. Fold until every piece is coated and no white patches remain.
    \end{enumerate}

    \item Pack the kimchi into the four 1-quart jars. Press each handful down hard so brine rises and air pockets disappear. Fill only to the shoulders, leaving about 1\nicefrac{1}{2}~inches of headspace. Press again after \textit{10~minutes} until brine covers the cabbage; the salt will keep drawing liquid. Trim each of the 4 reserved outer leaves to fit a jar mouth. Press one leaf on top of the kimchi, \emph{inside} the jar, so stray pieces stay under the brine. Screw on a plastic mason-jar storage lid (the one-piece cap sold for leftover storage, not the two-piece metal canning lid) until fingertip-tight. Salt and acid rust metal lids. Do not pack to the rim---the jars will overflow when they bubble.

    \item Leave the jars at cool room temperature, \textit{64--70°F}, out of sun, for \textit{24--48~hours}. The kimchi is ready to chill when you press the cabbage and bubbles rise through cloudy red brine, and a taste is lightly tangy, still crunchy, and smells sour-garlicky, not rotten. At \textit{72°F} or warmer, check at \textit{18--24~hours}. At \textit{64°F} or cooler, it may need \textit{2--3~days}. Burp a lid that domes, then close it again.

    \item Move the jars to the refrigerator at \textit{40°F} or below. Keep the lids closed so the remaining gas stays in the brine. Taste after \textit{3--4~days} cold. Peak sparkle is the first \textit{2~weeks}. Use a clean spoon every time, and keep the cabbage under the brine.
\end{enumerate}

\newpage
{\footnotesize
\setlength{\columnsep}{20pt}
\setlength{\multicolsep}{6pt}
\begin{multicols}{2}
\setlength{\parindent}{0pt}
\setlength{\parskip}{2pt}
\setlength{\itemsep}{0pt}
\setlength{\parsep}{0pt}

\subsection*{Yield}
\begin{itemize}
    \item Fills four 1-quart wide-mouth jars, packed to the shoulders (about 3--3\nicefrac{1}{2}~quarts)
    \item Banchan for many meals; open one jar at a time
\end{itemize}

\subsection*{Equipment Required}
\begin{itemize}
    \item 12-qt. mixing bowl
    \item Large colander
    \item Small saucepan and whisk
    \item Box grater
    \item Chef's knife and cutting board
    \item Measuring cups and spoons
    \item Four 1-quart wide-mouth mason jars
    \item Plastic mason-jar storage lids, wide-mouth (one-piece, not metal canning lids)
    \item Disposable gloves
    \item Optional: glass fermentation weights
\end{itemize}

\subsection*{Mise en Place}
\begin{itemize}
    \item Cut 5--6~lb. \textbf{napa cabbage} into 2~inch pieces; set aside 4 large outer leaves
    \item Mince \textbf{garlic} to \nicefrac{1}{2}~cup; grate 1~Tbsp. \textbf{ginger}, \nicefrac{1}{2}~medium \textbf{onion}, and \nicefrac{1}{2}~medium \textbf{Fuji apple}
    \item Cut 12~oz. \textbf{daikon} and 1 medium \textbf{carrot} into thin matchsticks; cut 8~\textbf{scallions} into 1~inch lengths
    \item Wash and dry the four 1-quart jars before packing
    \item Make the porridge while the cabbage salts so it has time to cool
\end{itemize}

\subsection*{Ingredient Tips}
\begin{itemize}
    \item \textbf{Gochugaru} is required. Cayenne, crushed red pepper, chili powder, and gochujang are the wrong flavor. Buy coarse Korean flakes (Tae-kyung, Wang, or Mother-in-Law's).
    \item \textbf{Fish sauce}: Red Boat, Three Crabs, or Squid brand. This recipe uses fish sauce only; salted shrimp is optional in other kimchi, not here.
    \item Morton kosher salt is denser than Diamond Crystal. Use 9~Tbsp. Morton in place of \nicefrac{3}{4}~cup Diamond Crystal, and still salt to the bend test.
    \item Skip iodized table salt for the brine if you can; iodine and anti-caking agents can slow fermentation.
    \item \textbf{Daikon} stands in for Korean radish. It is a little wetter; keep the matchsticks thin so they pickle with the cabbage.
    \item \textbf{Fuji apple} feeds the bacteria that make bubbles. Any sweet-crisp apple works. Korean pear is the traditional fruit if you find it.
    \item \textbf{Mochiko} is glutinous (sweet) rice flour. Regular rice flour will not thicken the same way. Last resort: cook 1~Tbsp. cornstarch in the 2~cups \textbf{water} until thick.
    \item Napa cabbage is sometimes labeled Chinese cabbage. Use the pale, crinkly heads, not green round cabbage.
\end{itemize}

\subsection*{Preparation Tips}
\begin{itemize}
    \item The bend test matters more than the clock. Underalting leaves the cabbage raw and prone to spoilage. Over-salting, if you skip a thorough rinse, stalls the fizz.
    \item Cool porridge completely. Warm paste makes the garlic harsh and the jars watery.
    \item The outer cabbage leaf is a cap, not a garnish. It sits on the kimchi inside the jar and holds loose pieces under the brine, where they will not mold. Discard it when you start eating that jar, or chop it in.
    \item Do not use the two-piece metal canning lid (flat plus ring) for storage. Salt and acid rust the metal. Buy wide-mouth plastic mason storage caps. They screw on like a jam lid. If you only have metal, use it for the counter ferment, then switch to plastic in the fridge.
    \item Fill jars only to the shoulders. Overflow is the usual mess in the first two days.
    \item Bubbles in the brine are the point. Do not throw the kimchi out for fizz. Discard for fuzzy mold (green, black, or blue), a rotten smell, or slimy cabbage. A thin white film on the brine can be skimmed; keep the cabbage submerged.
    \item A warm kitchen skips the bubbly stage and goes sour and soft. Aim for \textit{64--70°F} on the counter.
\end{itemize}

\subsection*{Make Ahead \& Storage}
\begin{itemize}
    \item Counter ferment \textit{24--48~hours} at \textit{64--70°F}, then refrigerate at \textit{40°F} or below.
    \item Best sparkle: \textit{3--14~days} in the fridge. Still crunchy and tangy through \textit{3--8~weeks}.
    \item Good to eat \textit{3--6~months} cold. After that it is sour and softer; use it for stew and fried rice.
    \item Open one jar at a time. Sealed jars keep their fizz longer.
    \item Burp a tight lid if it domes in the first \textit{2--3~weeks} cold, then close it again.
    \item Keep cabbage under brine. Use a clean spoon. Do not freeze.
\end{itemize}

\subsection*{Serving Suggestions}
\begin{itemize}
    \item Cold, as banchan, with hot rice.
    \item Alongside Korean BBQ drumsticks, grilled meats, or Quick Pickled Cucumbers.
    \item When it is very sour: kimchi jjigae, kimchi fried rice, or a grilled cheese.
    \item Peak eating is cold from the jar, while it still snaps and the brine still tingles.
\end{itemize}

\end{multicols}
}

\end{document}
